\documentclass[a4paper,12pt]{article}

\usepackage[utf8]{inputenc}
\usepackage[T1]{fontenc}
\usepackage[english]{babel}
\usepackage{amsmath}
\usepackage{amssymb}
\usepackage{graphicx}
\usepackage{hyperref}
\usepackage{geometry}
\usepackage{fontspec}
\usepackage{booktabs}
\usepackage{amsthm}
\usepackage{enumitem}
\usepackage{caption}


\geometry{margin=2.5cm}
\setmainfont{Times New Roman}

\geometry{a4paper, margin=1in}

\title{Cryptography Notes}
\author{Alessandro Vulcu}
\date{\today}

\newtheorem{theorem}{Theorem}[section]
\newtheorem{definition}{Definition}[section]
\newtheorem{proposition}{Proposition}[section]
\newtheorem{corollary}{Corollary}[section]
\newtheorem*{remark}{Remark}

\begin{document}

\maketitle
\newpage
\tableofcontents
\newpage

%\input{Lessons/Introduction.tex}
%\newpage

%\input{Lessons/Chapter2.tex}
%\newpage

%====================Lesson 2=========================%
\section{Divisibility}
\title{Technical Report on Divisibility, Euclidean Division, and Number Representation}

\subsection{Introduction to Divisibility}
Divisibility is a fundamental pillar of number theory, providing the essential language for describing the relationships between integers. This concept serves as the bedrock upon which more advanced topics are built, including the computation of greatest common divisors via the Euclidean algorithm, the unique factorization of integers into primes, and the representation of numbers in different bases. This section formally defines this core relationship and explores its most immediate and essential properties.

\begin{center}
\line(1,0){400}
\end{center}

\subsubsection*{Definition of Divisibility}
The relationship of divisibility is formally defined as follows:
Let $a, b$ be integers, with $b \neq 0$. We say that $b$ divides $a$ (written as $b \mid a$) if there exists an integer $c$ such that $a = bc$.
For instance, we can state that $2 \mid 4$ because $4 = 2 \times 2$. Conversely, $2$ does not divide $5$, as there is no integer $c$ for which $5 = 2c$.
It is critical to note that the definition explicitly prohibits division by zero; the divisor $b$ must be a non-zero integer. This leads to a foundational question: ``Does $1$ divide any integer?'' Based on the definition, the answer is yes. For any integer $a$, we can write $a = 1 \times a$, satisfying the condition with $c = a$. These elementary definitions give rise to several essential properties that govern the arithmetic of integers.

\subsection{Fundamental Properties and Proof}
The formal definition of divisibility leads directly to a set of properties that are instrumental in number theory.

\begin{proposition}
Let $a, b$, and $c$ be integers. The following properties hold:
\begin{enumerate}
    \item If $a \mid 1$, then $a = \pm 1$.
    \item If $a$ and $b$ are non-zero integers such that $a \mid b$ and $b \mid a$, then $a = \pm b$.
    \item If $a \mid b$ and $b \mid c$, then $a \mid c$ (Transitivity).
\end{enumerate}
\end{proposition}

\begin{center}
\line(1,0){400}
\end{center}

\begin{proof}
\begin{enumerate}
    \item If $a \mid 1$: By definition, there is an integer $c$ such that $1 = ac$. Since $a$ and $c$ are integers, the only possible integer solutions are $(a=1, c=1)$ or $(a=-1, c=-1)$. Therefore, $a = \pm 1$.
    \item If $a \mid b$ and $b \mid a$: By definition, there exist integers $c_1$ and $c_2$ such that $b = ac_1$ and $a = bc_2$. Substituting the first equation into the second gives $a = (ac_1)c_2 = a(c_1c_2)$. Since $a \neq 0$, we can divide by $a$ to get $1 = c_1c_2$. As $c_1$ and $c_2$ are integers, the only solutions are $(c_1=1, c_2=1)$ or $(c_1=-1, c_2=-1)$.
    \begin{itemize}
        \item If $c_1 = 1$, then $b = a \times 1 = a$.
        \item If $c_1 = -1$, then $b = a \times (-1) = -a$.
    \end{itemize}
    Thus, $a = \pm b$.
    \item If $a \mid b$ and $b \mid c$: By definition, there exist integers $c_1$ and $c_2$ such that $b = ac_1$ and $c = bc_2$. Substituting the first equation into the second gives $c = (ac_1)c_2 = a(c_1c_2)$. Since $c_1c_2$ is an integer, this satisfies the definition of divisibility, so $a \mid c$.
\end{enumerate}
\end{proof}

\begin{center}
\line(1,0){400}
\end{center}

This demonstrates the transitive nature of the divisibility relation. A particularly common instance of divisibility is divisibility by $2$, which gives rise to its own specific terminology.

\subsection{Even and Odd Integers}
An even number is an integer that is divisible by $2$; an integer that is not divisible by $2$ is an odd number. For example, $-4$ is even because $-4 = 2 \times (-2)$, whereas $13$ is odd.
By convention, when we refer to the ``divisors'' of a natural number, we are typically referring to its positive divisors. For example, while it is true that $-2 \mid 6$, the set of divisors of $6$ is conventionally listed as $\{1, 2, 3, 6\}$.
When an integer does not divide another perfectly, the concept of a remainder becomes crucial. This idea is formalized by the Euclidean division algorithm.

\section{The Euclidean Division Algorithm}

\subsection{The Division Theorem}
Euclidean division is a cornerstone of number theory that generalizes the concept of divisibility. It provides a formal and precise description of the outcome of integer division, asserting that dividing any integer $a$ by a positive integer $b$ yields a unique integer quotient $q$ and a unique integer remainder $r$.

\begin{proposition}[Euclidean Division]
Let $a$ and $b$ be natural numbers with $b > 0$. Then there exist unique integers $q$ (the quotient) and $r$ (the remainder) such that $a = bq + r$, with $0 \leq r < b$.
\end{proposition}
A direct and important consequence of this theorem is that $b$ divides $a$ if and only if the remainder $r = 0$.
For example, if $a = 27$ and $b = 7$, we can write $27 = 7 \times 3 + 6$. Here, the quotient $q$ is $3$ and the remainder $r$ is $6$, which satisfies the condition $0 \leq 6 < 7$. While the theorem guarantees the existence and uniqueness of $q$ and $r$, the next section details the practical method for their calculation.

\subsection{Practical Computation of Quotient and Remainder}
The quotient and remainder from Euclidean division can be computed directly using standard arithmetic operations. Starting with the identity $a/b = q + r/b$, we know from the theorem that $0 \leq r < b$, which implies $0 \leq r/b < 1$. This leads to the conclusion that the quotient $q$ is the integer part of the fraction $a/b$. This is formally expressed using the floor function:
$$ \text{Quotient: } q = \lfloor a/b \rfloor $$
Once the quotient is found, the remainder can be derived from the original equation $a = bq + r$:
$$ \text{Remainder: } r = a - bq = a - b \times \lfloor a/b \rfloor $$
To illustrate, let $a = 15476$ and $b = 137$.
\begin{itemize}
    \item Calculate the ratio: $15476 / 137 \approx 112.96\dots$
    \item The quotient is the floor of this value: $q = \lfloor 112.96\dots \rfloor = 112$.
    \item The remainder is then: $r = 15476 - 137 \times 112 = 15476 - 15344 = 132$.
\end{itemize}
The repeated application of this division process is the key to representing numbers in different bases, most notably in the binary system.

\section{Binary Representation of Natural Numbers}

\subsection{Definition and Process}
Binary representation is a system for expressing any natural number as a sum of powers of $2$, using only the digits $0$ and $1$. This base-$2$ system is foundational to all modern digital computing. The algorithm to convert a number to its binary form is a direct application of the principles of Euclidean division.

\subsubsection*{Definition}
Let $m \in \mathbb{N}$. We write $m = (m_{B-1}\dots m_1m_0)_2$ if $m = m_0 + m_1 \times 2 + \dots + m_{B-1}2^{B-1}$ and $m_i \in \{0, 1\}$. The values $m_i$ are known as bits.
The process for converting a number to binary involves repeatedly dividing the number by $2$ and recording the remainders. The remainders, read in reverse order of their calculation, form the binary representation.

\begin{center}
\captionof{table}{Conversion of Decimal 112 to Binary}
\begin{tabular}{|c|c|c|}
\hline
\textbf{Number ($m$)} & \textbf{Quotient ($q$)} & \textbf{Remainder ($r$) = Bit} \\
\hline
$112 \div 2$ & 56 & 0 ($m_0$) \\
\hline
$56 \div 2$ & 28 & 0 ($m_1$) \\
\hline
$28 \div 2$ & 14 & 0 ($m_2$) \\
\hline
$14 \div 2$ & 7 & 0 ($m_3$) \\
\hline
$7 \div 2$ & 3 & 1 ($m_4$) \\
\hline
$3 \div 2$ & 1 & 1 ($m_5$) \\
\hline
$1 \div 2$ & 0 & 1 ($m_6$) \\
\hline
\end{tabular}
\end{center}
Reading the remainders from bottom to top, we find that $112$ in decimal is $(1110000)_2$ in binary. This algorithmic process is guaranteed to terminate and provides a unique representation for any natural number, as established by the following theorem.

\subsection{Theorem on Binary Representation and Proof}

\begin{proposition}
Let $m \in \mathbb{N}$, with $m < 2^B$. The repeated division-by-2 process:
\begin{align*}
m &= 2q_0 + m_0 \\
q_0 &= 2q_1 + m_1 \\
q_1 &= 2q_2 + m_2 \\
\dots
\end{align*}
terminates with $q_{B-1} = 0$, and $m = m_0 + 2m_1 + \dots + 2^{B-1}m_{B-1}$.
\end{proposition}

\begin{center}
\line(1,0){400}
\end{center}

\begin{proof}
The proof consists of two parts: demonstrating the polynomial form and proving that the process terminates.

\subsubsection*{Polynomial Form}
We can show the relationship by algebraic substitution.
\begin{align*}
m &= 2q_0 + m_0 \\
\intertext{Substitute $q_0 = 2q_1 + m_1$:}
m &= 2(2q_1 + m_1) + m_0 = 2^2q_1 + 2m_1 + m_0 \\
\intertext{Substitute $q_1 = 2q_2 + m_2$:}
m &= 2^2(2q_2 + m_2) + 2m_1 + m_0 = 2^3q_2 + 2^2m_2 + 2m_1 + m_0
\end{align*}
Continuing this process until the final step where $q_{B-2} = 2q_{B-1} + m_{B-1}$, we arrive at the final form: $m = m_0 + 2m_1 + 2^2m_2 + \dots + 2^{B-1}m_{B-1}$.

\subsubsection*{Termination}
The process is guaranteed to terminate because the sequence of quotients is strictly decreasing and bounded.
\begin{itemize}
    \item Since $m < 2^B$, the first division $m = 2q_0 + m_0$ implies $2q_0 \leq m < 2^B$, so $q_0 < 2^{B-1}$.
    \item Similarly, $q_0 < 2^{B-1}$ implies $2q_1 \leq q_0 < 2^{B-1}$, so $q_1 < 2^{B-2}$.
\end{itemize}
This pattern continues: $q_i < 2^{B-1-i}$.
For the final quotient $q_{B-1}$, we have $q_{B-1} < 2^{B-1-(B-1)} = 2^0 = 1$. Since $q_{B-1}$ is a non-negative integer, the only possibility is $q_{B-1} = 0$. The process terminates when the quotient becomes zero.
\end{proof}

\begin{center}
\line(1,0){400}
\end{center}

This theorem also establishes a direct relationship between a number's magnitude and the number of bits required to represent it.

\subsection{Properties of Binary Representation}

\begin{remark}
If $m$ has $B$ bits in its binary representation, i.e., $m = (m_{B-1}\dots m_0)_2$, then $m < 2^B$.
\end{remark}

\begin{proof}
The largest possible value for a $B$-bit number occurs when all bits are $1$. This value can be calculated by summing the geometric series:
$$ m \leq 1 + 2 + 4 + \dots + 2^{B-1} = \frac{2^B - 1}{2 - 1} = 2^B - 1 $$
Therefore, $m < 2^B$.
\end{proof}

\begin{center}
\line(1,0){400}
\end{center}

This leads to a corollary that precisely characterizes the number of bits required for a given integer.

\begin{corollary}
A natural number $m$ has a binary representation with $B$ bits (where the most significant bit $m_{B-1} = 1$) if and only if $2^{B-1} \leq m < 2^B$.
\end{corollary}
This relationship provides a practical formula for finding the required number of bits, $B$, for any number $m$ using logarithms:
$$ B = \lfloor \log_2 m \rfloor + 1 $$
For example, to find the number of bits required to represent the number $368932$:
\begin{itemize}
    \item Calculate the base-2 logarithm: $\log_2(368932) \approx 18.4\dots$
    \item Apply the formula: $B = \lfloor 18.4\dots \rfloor + 1 = 18 + 1 = 19$. It requires $19$ bits.
\end{itemize}
The Euclidean division algorithm is not only useful for number representation but is also the most efficient method for finding the greatest common divisor of two integers.


\newpage
\section{The Greatest Common Divisor (GCD)}

\begin{definition}[Greatest Common Divisor]
Let $a, b \in \mathbb{Z}$, such that $a \neq 0$ or $b \neq 0$. The set of their common divisors has a largest element $d \in \mathbb{N}$, called the greatest common divisor of $a$ and $b$. We write $d = \gcd(a, b)$.
\end{definition}

\bigskip

\begin{remark}[Example of GCD]
To find $\gcd(18, 12)$, we can list the divisors of each number:
\begin{itemize}
    \item Divisors of 12: $\{1, 2, 3, 4, 6, 12\}$
    \item Divisors of 18: $\{1, 2, 3, 6, 9, 18\}$
\end{itemize}
The set of common divisors is $\{1, 2, 3, 6\}$. The largest element in this set is $6$, so $\gcd(18, 12) = 6$.
\par
While listing all divisors is feasible for small numbers, it is impractical for large integers. A far more powerful and efficient method is the Euclidean Algorithm.
\end{remark}


\newpage
\section{The Euclidean Algorithm}

\begin{theorem}[Euclidean Algorithm]
Let $a, b \in \mathbb{N} : a > 0, b > 0 \text{ and } a \geq b$. The following algorithm computes $\gcd(a, b)$ in a finite number of steps:
\begin{enumerate}
    \item Let $r_0 := a \text{ and } r_1 := b$.
    \item Set $i = 1$.
    \item Divide $r_{i-1}$ by $r_i$ to obtain a quotient $q_i$ and remainder $r_{i+1}$: $r_{i-1} = r_i q_i + r_{i+1}$.
    \item If $r_{i+1} = 0$, then $r_i = \gcd(a, b)$. The algorithm terminates.
    \item Otherwise, increment $i \rightarrow i+1$ and return to Step 3.
\end{enumerate}
\end{theorem}
\begin{proof}
Let the sequence of remainders be $r_0, r_1, \dots, r_n, r_{n+1}$ where $d = r_n$ is the last non-zero remainder and $r_{n+1} = 0$. We must demonstrate two facts: that $d$ is a common divisor of $a$ and $b$, and that it is the greatest common divisor.

\subsubsection*{a. $d \mid a$ and $d \mid b$}
We proceed by working backward from the end of the algorithm.
\begin{itemize}
    \item The final division step is $r_{n-1} = r_n q_n + r_{n+1}$. Since $r_{n+1} = 0$ and $d = r_n$, this equation becomes $r_{n-1} = d q_n$. This demonstrates $d \mid r_{n-1}$.
    \item The preceding step is $r_{n-2} = r_{n-1} q_{n-1} + r_n$. Since $d \mid r_n$ and we have shown $d \mid r_{n-1}$, it follows that $d$ must divide their linear combination, $d \mid r_{n-2}$.
\end{itemize}
Continuing this backward induction process, we establish that $d$ divides every preceding remainder. This concludes that $d \mid r_1$ (where $r_1 = b$) and $d \mid r_0$ (where $r_0 = a$). Thus, $d$ is a common divisor.

\subsubsection*{b. $d$ is the Greatest Common Divisor}
Let $c \in \mathbb{Z}$ be any common divisor of $a$ and $b$. We work forward down the chain of remainders.
\begin{itemize}
    \item Since $c \mid a$ ($r_0$) and $c \mid b$ ($r_1$), and $r_2 = r_0 - r_1q_1$, it must be that $c \mid r_2$.
    \item Since $c \mid r_1$ and $c \mid r_2$, and $r_3 = r_1 - r_2q_2$, it must be that $c \mid r_3$.
\end{itemize}
Continuing this process, we eventually find that $c$ must divide the last non-zero remainder, $r_n$, so $c \mid d$.
Since $c \mid d$, it is a consequence of divisibility that $c \leq d$. This holds $\forall c$ common divisor, proving that $d$ is the greatest common divisor. \qedhere
\end{proof}

\begin{center}
    \vspace{10pt}
    \line(1,0){450}
    \vspace{25pt}
\end{center}

\begin{proposition}[Extended Euclidean Algorithm]
Let $a, b \in \mathbb{N}$ such that $a \neq 0$ or $b \neq 0$. If $d = \gcd(a, b)$, then $\exists u, v \in \mathbb{Z} : d = au + bv$.
\end{proposition}

\begin{proof}
    The proof relies on the standard Euclidean Algorithm applied to $a$ and $b$, which generates a sequence of remainders $r_i$:
    \begin{align*}
        a &= q_1 b + r_1, \quad \text{with } 0 \le r_1 < b \\
        b &= q_2 r_1 + r_2, \quad \text{with } 0 \le r_2 < r_1 \\
        r_1 &= q_3 r_2 + r_3, \quad \text{with } 0 \le r_3 < r_2 \\
        &\vdots \\
        r_{n-3} &= q_{n-1} r_{n-2} + r_{n-1}, \quad \text{with } 0 \le r_{n-1} < r_{n-2} \\
        r_{n-2} &= q_n r_{n-1} + 0.
    \end{align*}
    
    Since $r_{n-1}$ is the last non-zero remainder, we have $d = \gcd(a, b) = r_{n-1}$.
    
    We proceed by **backward substitution** (as indicated by the notes on the right):
    \begin{itemize}
        \item From the second-to-last equation, $r_{n-1}$ can be expressed as a linear combination of $r_{n-3}$ and $r_{n-2}$:
        $$d = r_{n-1} = r_{n-3} - q_{n-1} r_{n-2}$$
        \item Since $r_{n-2}$ is a linear combination of $r_{n-4}$ and $r_{n-3}$, substituting this expression for $r_{n-2}$ into the equation for $d$ will express $d$ as a linear combination of $r_{n-4}$ and $r_{n-3}$.
        \item By continuing this backward substitution process (i.e., replacing each remainder $r_i$ with its expression from the previous division step in the Euclidean algorithm), we progressively express $d$ as a linear combination of the previous two terms in the remainder sequence:
        $$r_{i} = \text{lin. comb. of } r_{i-2} \text{ and } r_{i-1}$$
        \item This process continues until we reach the first two equations, where $d$ will be expressed as a linear combination of $a$ and $b$:
        $$d = r_{n-1} = au + bv$$
        for some integers $u, v \in \mathbb{Z}$.
    \end{itemize}
\end{proof}

\begin{center}
    \vspace{10pt}
    \line(1,0){450}
    \vspace{25pt}
\end{center}

\newpage
\begin{proposition}[Solvability of Linear Diophantine Equations]
    Let $a, b \in \mathbb{Z}$ be integers, not both zero. The equation $au + bv = c$ has integer solutions $u, v \in \mathbb{Z}$ if and only if $\gcd(a, b) \mid c$.
\end{proposition}

\begin{proof}
    \textbf{$(\impliedby)$ If $\gcd(a, b) \mid c$, solutions exist} \\
    Assume $c = k \cdot \gcd(a, b)$ for some $k \in \mathbb{Z}$. By Bézout's Identity, $\exists u_0, v_0 \in \mathbb{Z}$ so that $au_0 + bv_0 = \gcd(a, b)$.
    
    Multiplying by $k$ yields:
    $$k(au_0 + bv_0) = k \cdot \gcd(a, b) \implies a(ku_0) + b(kv_0) = c$$
    
    Thus, $u = ku_0$ and $v = kv_0$ are integer solutions.

    \textbf{$(\implies)$ If solutions exist, then $\gcd(a, b) \mid c$} \\
    Assume $au + bv = c$ has integer solutions $u, v \in \mathbb{Z}$.
    
    Let $d = \gcd(a, b)$. Since $d \mid a$ and $d \mid b$, $d$ must divide any linear combination of $a$ and $b$:
    $$d \mid (au + bv) \implies d \mid c$$
\end{proof}

\begin{center}
    \vspace{10pt}
    \line(1,0){450}
    \vspace{25pt}
\end{center}
\begin{proposition}[All Integer Solutions to $au+bv=c$]
    Let $a, b \in \mathbb{Z}$ be relatively prime (i.e., $\gcd(a, b) = 1$), and let $u_0, v_0 \in \mathbb{Z}$ be a particular solution so that $au_0 + bv_0 = c$. Then, all integer solutions $(u, v)$ to $au + bv = c$ are of the form:
    $$u = u_0 + kb \quad \text{and} \quad v = v_0 - ka$$
    for any $k \in \mathbb{Z}$.
\end{proposition}

\begin{proof}
    \textbf{Part 1: Verification (The form gives solutions)} \\
    Substituting the general form into the equation:
    \begin{align*}
        a(u_0 + kb) + b(v_0 - ka) &= au_0 + akb + bv_0 - bka \\
        &= (au_0 + bv_0) + (akb - bka) \\
        &= c + 0 \\
        &= c
    \end{align*}
    Thus, $\forall k \in \mathbb{Z}$, the formula yields a solution.

    \textbf{Part 2: Completeness (All solutions must be of this form)} \\
    Assume $(u, v)$ is any arbitrary solution, so $au + bv = c$. We also have the particular solution $au_0 + bv_0 = c$.
    
    Subtracting the two equations gives:
    $$a(u - u_0) + b(v - v_0) = 0 \implies a(u - u_0) = -b(v - v_0) \implies a(u - u_0) = b(v_0 - v)$$
    
    Since $a \mid b(v_0 - v)$ and we are given $\gcd(a, b) = 1$, by Euclid's Lemma, it must be that $a$ divides the other factor, $(v_0 - v)$:
    $$a \mid (v_0 - v)$$
    
    Let $k \in \mathbb{Z}$ be the integer so that $v_0 - v = ka$. Rearranging this gives the form for $v$:
    $$v = v_0 - ka$$
    
    Substituting $v_0 - v = ka$ back into the equality $a(u - u_0) = b(v_0 - v)$:
    $$a(u - u_0) = b(ka) = a(kb)$$
    
    Since $a \ne 0$, we can divide by $a$ to get:
    $$u - u_0 = kb$$
    
    Rearranging this gives the form for $u$:
    $$u = u_0 + kb$$
    
    Hence, all solutions are of the specified form.
\end{proof}

\begin{center}
    \line(1,0){450}
\end{center}

\newpage
\section{Coprime Integers}

\subsection{Definition and Characterization}
Coprime (or relatively prime) integers are those that share no common divisors other than $1$ and $-1$. This relationship is central to many theorems in number theory, including those related to prime numbers and factorization.

\subsubsection*{Definition}
Two integers $a, b$ are relatively prime (or coprime) if $\gcd(a, b) = 1$.
For example, $32$ and $15$ are relatively prime. The divisors of $32$ are $\{1, 2, 4, 8, 16, 32\}$ and the divisors of $15$ are $\{1, 3, 5, 15\}$. Their only common positive divisor is $1$.

\begin{proposition}[Characterisation of Coprime Numbers]
Integers $a$ and $b$ are relatively prime if and only if there exist integers $u, v$ such that $au + bv = 1$.
\end{proposition}

\begin{center}
\line(1,0){400}
\end{center}

\begin{proof}
\begin{itemize}
    \item[($\Rightarrow$)] This is a direct consequence of the Extended Euclidean Algorithm. If $\gcd(a, b) = 1$, then the algorithm guarantees the existence of integers $u, v$ such that $au + bv = 1$.
    \item[($\Leftarrow$)] Conversely, assume there are integers $u, v$ such that $au + bv = 1$. Let $d = \gcd(a, b)$. Since $d \mid a$ and $d \mid b$, it must divide their linear combination $au + bv$. Therefore, $d \mid 1$. The only positive integer that divides $1$ is $1$, so $d = 1$.
\end{itemize}
\end{proof}

\begin{center}
\line(1,0){400}
\end{center}

This characterization leads directly to an essential property regarding prime numbers.

\subsection{The Prime Divisor Property}

\begin{corollary}
Let $p$ be a prime number and $m$ be an integer. If $p$ does not divide $m$, then $p$ and $m$ are coprime.
\end{corollary}

\begin{proof}
The divisors of a prime number $p$ are only $\{1, p\}$. If $p$ does not divide $m$, then $p$ is not a common divisor. Therefore, the only common divisor of $p$ and $m$ is $1$, which means $\gcd(p, m) = 1$.
\end{proof}

\begin{center}
\line(1,0){400}
\end{center}

This corollary is the key to proving the following critical proposition, which is commonly known as Euclid's Lemma.

\begin{proposition}[Prime Divisor Property]
Let $p$ be a prime and $a, b$ be integers. If $p \mid ab$, then $p \mid a$ or $p \mid b$.
\end{proposition}

\begin{proof}
Assume $p \mid ab$. If $p \mid a$, the proposition is true. If $p$ does not divide $a$, then by the corollary above, $\gcd(p, a) = 1$. By the characterization of coprime numbers, there exist integers $u, v$ such that $pu + av = 1$. Multiplying this equation by $b$ gives $pub + abv = b$.
\begin{itemize}
    \item We know $p$ divides $pub$.
    \item We are given that $p$ divides $ab$, so it must also divide $abv$.
\end{itemize}
Since $p$ divides both terms on the left side, it must divide their sum. Therefore, $p \mid b$.
\end{proof}

\begin{center}
\line(1,0){400}
\end{center}

This property is unique to prime numbers. For a non-prime (composite) number, this does not hold. For example, $4 \mid 2 \times 6$ (since $4 \mid 12$), but $4$ does not divide $2$ and $4$ does not divide $6$. This unique property of primes is the cornerstone of the most important theorem in elementary number theory.

\section{Conclusion: The Fundamental Theorem of Arithmetic}
This report has traced a logical path from the elementary definition of divisibility to the profound and unique properties of prime numbers. We began by establishing the rules of integer division, which we generalized with the Euclidean Division theorem. This theorem, in turn, provided a practical algorithm not only for finding quotients and remainders---enabling number system conversions like binary representation---but also for efficiently computing the greatest common divisor. By extending the Euclidean algorithm, we established that the GCD can be expressed as a linear combination, a result that unlocked the solution to linear Diophantine equations and provided a powerful characterization for coprime integers.
The culmination of this journey is the Prime Divisor Property, which states that if a prime divides a product, it must divide one of the factors. This property, which we proved does not hold for composite numbers, is the essential lemma required to prove the Fundamental Theorem of Arithmetic. This theorem asserts that every integer greater than $1$ is either a prime number itself or can be represented as a unique product of prime numbers, cementing the role of primes as the fundamental building blocks of the integers.


%====================Lesson 3=========================%

\newpage
\section{Classes Modulo \texorpdfstring{$m$}{m}}

\begin{proposition}
\label{prop:euclidean-division-unique-remainder}
$\forall a \in \mathbb{Z}$ (the dividend) and $\forall m \in \mathbb{N}$ so that $m > 1$ (the modulus), $\exists! q, r \in \mathbb{Z} : a = mq + r$, with $0 \leq r < m$.
\end{proposition}

\begin{center}
    \line(1,0){450}
\end{center}

\begin{theorem}[Existence of a Unique Remainder]
$\forall a \in \mathbb{Z}$ and $\forall m \in \mathbb{N} : m > 1$, $\exists! r \in \{0, 1, \dots, m-1\} : a \equiv r \pmod{m}$.
\end{theorem}
\begin{proof}
By Proposition \ref{prop:euclidean-division-unique-remainder} (Euclidean Division), we have $\exists! q, r \in \mathbb{Z} : a = mq + r$, where $0 \leq r < m$.
Rearranging this equation gives $a - r = mq$.
By the formal definition of divisibility, $m \mid (a - r)$.
By the formal definition of congruence, this directly implies $a \equiv r \pmod{m}$. The uniqueness of $r$ is guaranteed by the Euclidean Division Theorem. \qedhere
\end{proof}

\begin{center}
    \line(1,0){450}
\end{center}

\begin{theorem}[Congruence is an Equivalence Relation]
The relation of congruence modulo $m$ (denoted by $\equiv \pmod{m}$) is an equivalence relation on the set of integers $\mathbb{Z}$.
\end{theorem}
\begin{proof}
We must demonstrate that congruence modulo $m$ satisfies reflexivity, symmetry, and transitivity.
\begin{itemize}
    \item \textbf{Reflexivity:} $\forall a \in \mathbb{Z}$, $a - a = 0$. Since $m \mid 0$, $a \equiv a \pmod{m}$. The relation is reflexive.
    \item \textbf{Symmetry:} Assume $a \equiv b \pmod{m}$. $\implies m \mid (a - b)$, so $\exists k \in \mathbb{Z} : a - b = k m$.
    Multiplying by $-1$ yields $b - a = (-k) m$. Since $-k \in \mathbb{Z}$, $\implies m \mid (b - a)$, so $b \equiv a \pmod{m}$. The relation is symmetric.
    \item \textbf{Transitivity:} Assume $a \equiv b \pmod{m}$ and $b \equiv c \pmod{m}$.
    $\implies \exists k_1, k_2 \in \mathbb{Z} : a - b = k_1 m \text{ and } b - c = k_2 m$.
    Adding the equations:
    $$ (a - b) + (b - c) = k_1 m + k_2 m $$
    $$ a - c = (k_1 + k_2) m $$
    Since $k_1 + k_2 \in \mathbb{Z}$, $\implies m \mid (a - c)$, so $a \equiv c \pmod{m}$. The relation is transitive. \qedhere
\end{itemize}
\end{proof}

\begin{center}
    \line(1,0){450}
\end{center}

\begin{proposition}[Equivalence Classes Form a Partition]
Let $\sim$ be an equivalence relation on a set $X$. The set of equivalence classes, $X/\sim$, forms a partition of $X$.
\end{proposition}

\begin{center}
    \line(1,0){450}
\end{center}

\newpage
\section{Invertibility Modulo \texorpdfstring{$m$}{m}}

\begin{proposition}[Compatibility]
Let $a_1, a_2, b_1, b_2 \in \mathbb{Z}$. If $a_1 \equiv a_2 \pmod{m}$ and $b_1 \equiv b_2 \pmod{m}$, then:
\begin{enumerate}
    \item $a_1 + b_1 \equiv a_2 + b_2 \pmod{m}$
    \item $a_1b_1 \equiv a_2b_2 \pmod{m}$
\end{enumerate}
\end{proposition}
\begin{proof}
    \leavevmode
\begin{itemize}
    \item \textbf{Addition:} Given $a_2 \equiv a_1 \pmod{m}$ and $b_2 \equiv b_1 \pmod{m}$, $\implies \exists k_1, k_2 \in \mathbb{Z} : a_2 - a_1 = k_1 m \text{ and } b_2 - b_1 = k_2 m$.
    Adding the equations yields $(a_2 - a_1) + (b_2 - b_1) = (k_1 + k_2) m$.
    Rearranging the terms: $(a_2 + b_2) - (a_1 + b_1) = (k_1 + k_2) m$.
    Since $k_1 + k_2 \in \mathbb{Z}$, this shows $m \mid ((a_2 + b_2) - (a_1 + b_1))$, so $a_1 + b_1 \equiv a_2 + b_2 \pmod{m}$.
    \item \textbf{Multiplication:} We express $a_2 = a_1 + k_1 m$ and $b_2 = b_1 + k_2 m$.
    Multiplying them gives:
    $$ a_2b_2 = (a_1 + k_1 m)(b_1 + k_2 m) = a_1b_1 + a_1 k_2 m + b_1 k_1 m + k_1 k_2 m^2 $$
    Factoring out $m$: $a_2b_2 - a_1b_1 = m(a_1 k_2 + b_1 k_1 + k_1 k_2 m)$.
    Let $K = a_1 k_2 + b_1 k_1 + k_1 k_2 m$. Since $K \in \mathbb{Z}$, $m \mid (a_2b_2 - a_1b_1)$, so $a_1b_1 \equiv a_2b_2 \pmod{m}$. \qedhere
\end{itemize}
\end{proof}

\begin{center}
    \vspace{10pt}
    \line(1,0){450}
    \vspace{25pt}
\end{center}

\begin{proposition}
The set of integers modulo $m$ under addition, $(\mathbb{Z}/m\mathbb{Z}, +)$, forms a \textbf{commutative group}.
\end{proposition}

\begin{center}
    \vspace{10pt}
    \line(1,0){450}
    \vspace{25pt}
\end{center}

\begin{proposition}
When considering both addition and multiplication, $(\mathbb{Z}/m\mathbb{Z}, +, \cdot)$ forms a \textbf{commutative ring}.
\end{proposition}

\begin{center}
    \vspace{10pt}
    \line(1,0){450}
    \vspace{25pt}
\end{center}

\begin{theorem}
If an inverse exists, it is unique.
\end{theorem}
\begin{proof}
Assume an element $[a]$ has two inverses, $[a']$ and $[a'']$. By definition, $[a] \cdot [a'] = [1]$ and $[a] \cdot [a''] = [1]$. Since multiplication is associative and $[1]$ is the identity element, we can write:
$$ [a'] = [a'] \cdot [1] = [a'] \cdot ([a] \cdot [a'']) $$
$$ [a'] = ([a'] \cdot [a]) \cdot [a''] = [1] \cdot [a''] = [a''] $$
Therefore, $[a'] = [a'']$, $\implies$ the inverse is unique.
\end{proof}

\begin{center}
    \vspace{10pt}
    \line(1,0){450}
    \vspace{25pt}
\end{center}

\begin{theorem}[Characterization of Invertibility]
An element $[a] \in \mathbb{Z}/m\mathbb{Z}$ is invertible $\iff \gcd(a, m) = 1$.
\end{theorem}
\begin{proof}
An element $[a]$ is invertible $\iff \exists u \in \mathbb{Z} : a u \equiv 1 \pmod{m}$.
This congruence is equivalent to the linear Diophantine equation $a u - 1 = -m v$ for some $v \in \mathbb{Z}$, which can be rewritten as $a u + m v = 1$.
By Bézout's Identity, this equation has an integer solution $\iff \gcd(a, m) \mid 1$.
Since $\gcd(a, m) \in \mathbb{N}$, this condition is satisfied $\iff \gcd(a, m) = 1$. \qedhere
\end{proof}

\begin{center}
    \vspace{10pt}
    \line(1,0){450}
    \vspace{25pt}
\end{center}

\begin{proposition}
The set of units under multiplication, $(\mathbb{Z}/m\mathbb{Z})^*$, forms a \textbf{commutative group} (the group of units modulo $m$).
\end{proposition}

\begin{center}
    \vspace{10pt}
    \line(1,0){450}
    \vspace{25pt}
\end{center}

\begin{proposition}
$\forall p$ prime number, the commutative ring $(\mathbb{Z}/p\mathbb{Z}, +, \cdot)$ is a \textbf{field}, commonly denoted $\mathbb{F}_p$.
\end{proposition}

\begin{center}
    \vspace{10pt}
    \line(1,0){450}
    \vspace{25pt}
\end{center}

\begin{corollary}
Let $p \in \mathbb{N}$ be a prime so that $p > 2$. The equation $x^2 \equiv 1 \pmod{p}$ in $\mathbb{F}_p$ has exactly two distinct solutions: $x \equiv 1 \pmod{p}$ and $x \equiv p-1 \pmod{p}$ (i.e., $x \equiv -1 \pmod{p}$).
\end{corollary}
\begin{proof}
The equation $x^2 \equiv 1 \pmod{p}$ is equivalent to $p \mid (x^2 - 1)$, or $p \mid (x - 1)(x + 1)$.
Since $p$ is prime, by the Prime Divisor Property:
\begin{itemize}
    \item Either $p \mid (x - 1)$, which implies $x \equiv 1 \pmod{p}$.
    \item Or $p \mid (x + 1)$, which implies $x \equiv -1 \equiv p-1 \pmod{p}$.
\end{itemize}
These two solutions are distinct because $1 \not\equiv -1 \pmod{p}$ for $p > 2$. \qedhere
\end{proof}

\begin{center}
    \vspace{10pt}
    \line(1,0){450}
    \vspace{25pt}
\end{center}

\newpage
\section{Application I: Solving Linear Congruences}
The algebraic theory provides a framework for solving linear congruences of the form $ax \equiv b \pmod{m}$.

\subsection{The General Condition for Solutions}
\begin{proposition}
The linear congruence $ax \equiv b \pmod{m}$ has at least one integer solution for $x$ if and only if $d = \gcd(a, m)$ divides $b$.
\end{proposition}
\begin{proof}
The congruence is equivalent to the linear Diophantine equation $ax + mv = b$. This equation is solvable if and only if $\gcd(a, m)$ divides $b$.
\end{proof}

\subsection{Case 1: Unique Solution ($\gcd(a,m) = 1$)}
When $\gcd(a, m) = 1$, the unique solution modulo $m$ is found by multiplying by the inverse $a^{-1}$:
$$ x \equiv a^{-1}b \pmod{m} $$
\begin{itemize}
    \item \textbf{Example:} To solve $245x \equiv 24 \pmod{677}$, since $677$ is prime, $245^{-1} \equiv 572 \pmod{677}$.
    $$ x \equiv 572 \cdot 24 \equiv 13728 \equiv 188 \pmod{677} $$
\end{itemize}

\subsection{Case 2: Multiple or No Solutions ($\gcd(a,m) = d > 1$)}
If $d = \gcd(a, m)$ does not divide $b$, there are no solutions.
\begin{proposition}
If $d \mid b$, the congruence $ax \equiv b \pmod{m}$ has exactly $d$ distinct solutions modulo $m$.
\end{proposition}
The method is:
\begin{enumerate}
    \item \textbf{Reduce:} Divide by $d$: $\frac{a}{d}x \equiv \frac{b}{d} \pmod{\frac{m}{d}}$.
    \item \textbf{Solve:} Find the unique solution $x_0$ to the reduced congruence.
    \item \textbf{Generate Solutions:} The complete set of $d$ solutions is:
    $$ x_k = x_0 + k\left(\frac{m}{d}\right) \quad \text{for } k = 0, 1, \dots, d-1 $$
\end{enumerate}
\begin{itemize}
    \item \textbf{Example: $9x \equiv 6 \pmod{15}$.}
    \begin{enumerate}
        \item $d = \gcd(9, 15) = 3$. Since $3 \mid 6$, there are 3 solutions.
        \item \textbf{Reduced Congruence:} $3x \equiv 2 \pmod{5}$. Solution $x_0 = 4$.
        \item \textbf{Solutions:} $m/d = 5$. Solutions are $x_0=4$, $x_1=4+5=9$, $x_2=4+10=14$.
    \end{enumerate}
\end{itemize}

\section{Application II: Shamir's Secret Sharing Scheme}
Shamir's Secret Sharing is an elegant application of modular arithmetic and linear algebra for secure information distribution.

\subsection{The Premise and Analogy}
The scheme relies on the principle that $t$ distinct points uniquely determine a polynomial of degree $t-1$. The secret is encoded as a coefficient (e.g., the y-intercept $a_0$) of a polynomial $P(x)$, and each person receives a point $(x_i, P(x_i))$ on the curve. All calculations are performed in a finite field $\mathbb{F}_p$.

\subsection{A Worked Example}
\begin{itemize}
    \item \textbf{Information:} Friend A: $(23, 401)$; Friend B: $(58, 368)$.
    \item \textbf{Public Parameter:} Modulus $p = 941$.
    \item \textbf{Task:} Calculate the secret slope $m$ (the secret) of the line passing through the points: $m = (y_2 - y_1) \cdot (x_2 - x_1)^{-1} \pmod{941}$.
\end{itemize}
\subsubsection*{Calculation}
\begin{enumerate}
    \item \textbf{Differences:} $y_2 - y_1 = -33$; $x_2 - x_1 = 35$.
    \item \textbf{Inverse:} Using the Extended Euclidean Algorithm, $35^{-1} \equiv 242 \pmod{941}$.
    \item \textbf{Solve:} $m \equiv (-33) \cdot 242 \equiv -7986 \pmod{941}$.
    \item \textbf{Reduce:} $-7986 = -9 \times 941 + 483$.
    $$ m \equiv 483 \pmod{941} $$
\end{enumerate}
The calculated secret slope is $m = 483$.

\newpage
\section{An Exposition on the Chinese Remainder Theorem}

\section{Introduction: An Ancient Problem with Modern Relevance}
The \textbf{Chinese Remainder Theorem (CRT)} is a cornerstone of number theory with a history as rich as its modern applications are diverse. Its origins can be traced to a problem of enumeration found in the 3rd to 5th century Chinese mathematical text, \textit{Sun Tze Suan Ching}. The text poses a now-famous puzzle:
\begin{quote}
If we count them by threes we have two left over. If we count them by fives we have three left over. If we count them by sevens we have two left over.
\end{quote}
This ancient riddle of remainders lays the groundwork for a profound mathematical principle. Centuries later, this principle has found new life as a foundational tool in modern applied mathematics, most notably in the field of cryptography. The theorem provides a powerful framework for solving systems of simultaneous congruences, a task central to the design and implementation of various cryptographic algorithms. This document provides a comprehensive, formal, and academic exploration of the Chinese Remainder Theorem, detailing its formal statement, rigorous proofs, practical solution methodologies, and critical extensions.

\section{The System of Congruences: Formalizing the Problem}
The first step in transforming an abstract puzzle into a solvable mathematical problem is to translate its language into a formal system. This act of formalization is crucial, as it allows us to move beyond a specific case and develop a general, repeatable methodology for all problems of a similar structure.
The classic problem described by Sun Tze can be stated as follows:
\begin{itemize}
    \item If we count them by threes we have two left over.
    \item If we count them by fives we have three left over.
    \item If we count them by sevens we have two left over.
\end{itemize}
In the language of modular arithmetic, this puzzle translates directly into a system of three linear congruences with a single unknown variable, $x$:
\begin{align*}
x &\equiv 2 \pmod{3} \\
x &\equiv 3 \pmod{5} \\
x &\equiv 2 \pmod{7}
\end{align*}
This formal system is the precise mathematical expression of the original problem. The Chinese Remainder Theorem provides the theoretical framework that guarantees a solution to such a system exists and offers a clear path to finding it.

\section{The Chinese Remainder Theorem: A Formal Statement}
The Chinese Remainder Theorem offers a powerful and elegant guarantee regarding the solution of systems of linear congruences. It asserts that under a specific but critical condition---that the moduli are pairwise relatively prime---a solution not only exists but is also unique within a well-defined modular range.

\begin{theorem}[Chinese Remainder Theorem]
Let $m_1, m_2, \dots, m_k$ be \textbf{pairwise relatively prime} integers. Let $a_1, a_2, \dots, a_k$ be arbitrary integers.
The system of congruences:
$$ x \equiv a_i \pmod{m_i} \quad \text{for } i = 1 \dots k $$
has a solution, and the complete set of solutions is congruent modulo the product $M = m_1 \cdot m_2 \cdot \dots \cdot m_k$.
\end{theorem}
The theorem makes two fundamental promises. The first is \textbf{existence}: a solution to the system is guaranteed to exist. The second is \textbf{uniqueness}: while there is an infinite set of solutions, they are all congruent to one another modulo $M$. This means there is exactly one solution within the range from $0$ to $M-1$.

\section{Solution Methodologies}
While the theorem guarantees a solution, several practical methods can be employed to find it. These approaches range from an intuitive, iterative process suitable for smaller systems to a more systematic and scalable algorithm derived directly from the theorem's constructive proof.

\subsection{Method 1: Iterative Substitution (``By Hands'')}
This intuitive method, suitable for smaller systems, involves solving the first congruence, substituting the resulting expression into the next congruence, and systematically solving for the introduced parameters.

\subsubsection*{Example 1: A 2x2 System}
Consider the system:
\begin{align*}
x &\equiv 3 \pmod{11} \\
x &\equiv 2 \pmod{7}
\end{align*}
\begin{enumerate}
    \item \textbf{Express $x$:} $x = 3 + 11k$.
    \item \textbf{Substitute:} $3 + 11k \equiv 2 \pmod{7}$.
    \item \textbf{Solve for $k$:} $4k \equiv -1 \equiv 6 \pmod{7}$. The inverse of $4$ modulo $7$ is $2$. Multiplying by $2$ yields $k \equiv 2 \cdot 6 \equiv 12 \equiv 5 \pmod{7}$.
    \item \textbf{Express $k$:} $k = 5 + 7L$.
    \item \textbf{Substitute back:} $x = 3 + 11(5 + 7L) = 58 + 77L$.
\end{enumerate}
Thus, the solution is $x \equiv 58 \pmod{77}$.

\subsubsection*{Example 2: The Sun Tze Problem}
\begin{align*}
x &\equiv 2 \pmod{3} \\
x &\equiv 3 \pmod{5} \\
x &\equiv 2 \pmod{7}
\end{align*}
\begin{enumerate}
    \item $x = 2 + 3k$. Substitute into second: $3k \equiv 1 \pmod{5}$. Inverse of 3 is 2, so $k \equiv 2 \pmod{5}$.
    \item $k = 2 + 5m$. Substitute back into $x$: $x = 2 + 3(2 + 5m) = 8 + 15m$.
    \item Substitute into third: $8 + 15m \equiv 2 \pmod{7}$. Since $15 \equiv 1$ and $8 \equiv 1 \pmod{7}$, we get $1 + m \equiv 2 \implies m \equiv 1 \pmod{7}$.
    \item $m = 1 + 7L$. Substitute back: $x = 8 + 15(1 + 7L) = 23 + 105L$.
\end{enumerate}
The smallest positive solution is $x = 23$.

\subsection{Method 2: The Constructive Algorithm}
This algorithm, derived from the theorem's constructive proof, is more robust and scalable for larger systems.
\begin{enumerate}
    \item \textbf{Step 1:} Define the moduli $m_i$ and calculate their product $M = m_1 \cdot \dots \cdot m_k$.
    \item \textbf{Step 2:} For each $i$, calculate $N_i = M / m_i$.
    \item \textbf{Step 3:} Find the modular multiplicative inverse $v_i$ of $N_i$ modulo $m_i$, such that $v_iN_i \equiv 1 \pmod{m_i}$. This is guaranteed by $\gcd(N_i, m_i) = 1$.
    \item \textbf{Step 4 (Solution):} The final solution $x$ is constructed by the sum:
    $$ x = \sum_{i=1}^{k} a_iv_iN_i $$
\end{enumerate}

\subsubsection*{Worked Example}
Solve the system:
\begin{align*}
x &\equiv 2 \pmod{3} \\
x &\equiv 3 \pmod{7} \\
x &\equiv 4 \pmod{16}
\end{align*}
$a_i: 2, 3, 4$. $m_i: 3, 7, 16$. $M = 3 \cdot 7 \cdot 16 = 336$.
\begin{enumerate}
    \item \textbf{Calculate $N_i$:} $N_1 = 112$, $N_2 = 48$, $N_3 = 21$.
    \item \textbf{Find Inverses $v_i$:}
    \begin{itemize}
        \item $112v_1 \equiv 1 \pmod{3} \implies 1v_1 \equiv 1 \implies v_1 = 1$.
        \item $48v_2 \equiv 1 \pmod{7} \implies 6v_2 \equiv 1 \implies v_2 = -1 \equiv 6$.
        \item $21v_3 \equiv 1 \pmod{16} \implies 5v_3 \equiv 1 \implies v_3 = -3 \equiv 13$.
    \end{itemize}
    \item \textbf{Construct Solution:}
    \begin{align*}
    x &= 2(1)(112) + 3(6)(48) + 4(13)(21) \\
    x &= 224 + 864 + 1092 = 2180
    \end{align*}
    \item \textbf{Reduce:} $x \equiv 2180 \equiv 164 \pmod{336}$.
\end{enumerate}
The smallest positive solution is $x = 164$.

\section{Formal Proof of the Chinese Remainder Theorem}
The proof is divided into two parts: Existence (constructive) and Uniqueness (based on coprimality).

\subsection{Proof of Existence}
We must show that $x = \sum a_j v_j N_j$ is a solution. We verify that $x \equiv a_i \pmod{m_i}$ for an arbitrary $i$.
The argument hinges on the fact that for any $j \neq i$, the term $N_j$ contains $m_i$ as a factor, leading to:
$$ a_j v_j N_j \equiv 0 \pmod{m_i} \quad \text{for } j \neq i $$
Thus, when $x$ is taken modulo $m_i$, all terms vanish except the $i$-th term:
$$ x \equiv a_i v_i N_i \pmod{m_i} $$
Since $v_i N_i \equiv 1 \pmod{m_i}$ by construction, the proof concludes:
$$ x \equiv a_i(1) \equiv a_i \pmod{m_i} $$
This confirms the existence of a solution.

\subsection{Proof of Uniqueness}
\begin{enumerate}
    \item \textbf{Assume Two Solutions:} Let $x$ and $x'$ both solve the system.
    \item \textbf{Relationship:} For every $i$, $x \equiv a_i \pmod{m_i}$ and $x' \equiv a_i \pmod{m_i}$, implying $x - x' \equiv 0 \pmod{m_i}$. Thus, $m_i$ divides $(x - x')$.
    \item \textbf{Use Co-prime Condition:} Since all $m_i$ are pairwise co-prime and each divides $(x - x')$, their product $M = m_1 \cdot \dots \cdot m_k$ must also divide $(x - x')$.
    \item \textbf{Conclusion:} By definition, $M \mid (x - x')$ means $x \equiv x' \pmod{M}$.
\end{enumerate}
This confirms that the solution is unique modulo $M$.

\section{An Important Extension: The Case of Non-Co-prime Moduli}
The powerful guarantees of the CRT rely entirely on the condition that the moduli are pairwise co-prime. When this fails, a solution is not guaranteed.

\subsubsection*{Example of No Solution}
The system $x \equiv 0 \pmod{4}$ and $x \equiv 1 \pmod{6}$ has no solution because $x$ must be both even (from $x \equiv 0 \pmod{4}$) and odd (from $x \equiv 1 \pmod{6}$), a contradiction.

\subsubsection*{Compatibility Condition}
For a system of congruences with non-co-prime moduli, a solution exists \textbf{if and only if} the following compatibility condition is met for every pair of equations:
$$ a_i \equiv a_j \pmod{\gcd(m_i, m_j)} \quad \text{for every pair } i, j $$
If this condition holds, a unique solution exists modulo the \textbf{least common multiple} (LCM) of the moduli.

\section{Conclusion}
The Chinese Remainder Theorem is a profound result in number theory that provides a powerful bridge between modular arithmetic with large composite numbers and systems of simpler congruences. Its core function is to transform a complex problem into a manageable set of smaller, independent problems. This ``divide and conquer'' approach, rooted in an ancient counting puzzle, has become an indispensable tool in modern fields. In cryptography, for example, algorithms like RSA leverage the CRT to dramatically speed up computations involving large moduli. By breaking a single, computationally expensive exponentiation into several smaller ones that can be executed in parallel, the CRT demonstrates its enduring relevance, transforming an elegant piece of ancient mathematics into a cornerstone of modern computational security.

\newpage
\section{Euler's \texorpdfstring{$\phi$}{phi} Function}
\begin{definition}[Euler's Totient Function, $\phi$]
    Let $m \in \mathbb{N}$ such that $m \ge 1$. The **Euler's totient function**, denoted by $\phi(m)$, is the number of elements in the multiplicative group of integers modulo $m$, which is $\left(\mathbb{Z}/m\mathbb{Z}\right)^*$.
    
    In other words, $\phi(m)$ is the number of integers $k$ such that $1 \le k \le m$ so that $\gcd(k, m) = 1$.
\end{definition}

\begin{center}
    \vspace{10pt}
    \line(1,0){450}
    \vspace{25pt}
\end{center}

\begin{definition}[Remainder Map]
    Let $m, n \in \mathbb{N}$ be integers. The **remainder map** $r$ is a function:
    $$r: \mathbb{Z}/(mn)\mathbb{Z} \longrightarrow (\mathbb{Z}/m\mathbb{Z}) \times (\mathbb{Z}/n\mathbb{Z})$$
    
    It is defined $\forall a \in \{0, \dots, mn - 1\}$ as:
    $$r([a]_{mn}) := ([a]_m, [a]_n)$$
    where $[a]_k$ denotes the residue class of $a$ modulo $k$.
\end{definition}

\begin{center}
    \vspace{10pt}
    \line(1,0){450}
    \vspace{25pt}
\end{center}

\bigskip

\begin{proposition}[Remainder Map is Well-Defined]
    The remainder map $r$ is well-defined. That is, the result of the map is independent of the choice of the representative of the residue class:
    $$[a]_{mn} = [b]_{mn} \implies r([a]_{mn}) = r([b]_{mn})$$
    which means $([a]_m, [a]_n) = ([b]_m, [b]_n)$.
\end{proposition}

\begin{proof}
    We assume the antecedent: $[a]_{mn} = [b]_{mn}$.
    
    By the definition of equality of residue classes, this means $mn$ divides the difference $(b - a)$:
    $$[a]_{mn} = [b]_{mn} \implies mn \mid (b - a)$$
    
    Since $mn \mid (b - a)$, it must be that both $m$ and $n$ individually divide the difference $(b - a)$:
    $$mn \mid (b - a) \implies \begin{cases} m \mid (b - a) \\ n \mid (b - a) \end{cases}$$
    
    By the definition of congruence and residue classes, this translates back into the desired equalities:
    $$\begin{cases} m \mid (b - a) \implies [a]_m = [b]_m \\ n \mid (b - a) \implies [a]_n = [b]_n \end{cases}$$
    
    Since both components are equal, the ordered pairs are equal:
    $$([a]_m, [a]_n) = ([b]_m, [b]_n)$$
    
    Thus, $r([a]_{mn}) = r([b]_{mn})$, and the map is well-defined.
\end{proof}

\begin{center}
    \vspace{10pt}
    \line(1,0){450}
    \vspace{25pt}
\end{center}


\begin{theorem}[Bijectivity of the Remainder Map with Coprime Factors]
    Let $m, n \in \mathbb{N}$ be integers such that $m, n \ge 1$, and assume $m$ and $n$ are **pairwise coprime** (i.e., $\gcd(m, n) = 1$).
    
    The remainder map $r$:
    $$r: \mathbb{Z}/(mn)\mathbb{Z} \longrightarrow (\mathbb{Z}/m\mathbb{Z}) \times (\mathbb{Z}/n\mathbb{Z})$$
    is a **bijection** (one-to-one and onto).
\end{theorem}

\begin{proof}
    The proof establishes bijectivity by demonstrating both **Surjectivity** (the map is onto) and **Injectivity** (the map is one-to-one).

    \subsection*{Surjectivity (Onto)}
    Let $([a]_m, [b]_n) \in (\mathbb{Z}/m\mathbb{Z}) \times (\mathbb{Z}/n\mathbb{Z})$ be an arbitrary element in the codomain.
    
    We seek to show $\exists x \in \mathbb{Z}$ so that $r([x]_{mn}) = ([a]_m, [b]_n)$.
    
    This is equivalent to finding an integer $x$ that simultaneously solves the system of congruences:
    $$\begin{cases} x \equiv a \pmod{m} \\ x \equiv b \pmod{n} \end{cases}$$
    
    Since we are given that $\gcd(m, n) = 1$, the **Chinese Remainder Theorem** guarantees that a solution $x$ always exists for any pair of residues $(a, b)$.
    
    Therefore, $\forall$ element in the codomain, there is a corresponding pre-image in the domain, proving that the map $r$ is **surjective**.
    
    \subsection*{Injectivity (One-to-One)}
    We can establish injectivity by counting the elements (since the sets are finite):
    
    \begin{enumerate}
        \item **Count the Domain:** The number of elements in the domain is the modulus:
        $$|\mathbb{Z}/(mn)\mathbb{Z}| = mn$$
        
        \item **Count the Codomain:** The number of elements in the codomain (a Cartesian product) is the product of the sizes of the factors:
        $$|(\mathbb{Z}/m\mathbb{Z}) \times (\mathbb{Z}/n\mathbb{Z})| = |\mathbb{Z}/m\mathbb{Z}| \cdot |\mathbb{Z}/n\mathbb{Z}| = m \cdot n = mn$$
        
        \item **Conclusion:** Since the domain and codomain are finite sets and have the **same number of elements** ($mn = mn$), and we have already proven the map $r$ is **surjective**, it must automatically be **injective** as well.
        
        $$\text{Surjectivity} \land (|\text{Domain}| = |\text{Codomain}|) \implies \text{Injectivity}$$
    \end{enumerate}
    
    Since $r$ is both surjective and injective, it is a **bijection**.
\end{proof}

\begin{center}
    \vspace{10pt}
    \line(1,0){450}
    \vspace{25pt}
\end{center}


\begin{theorem}[Totient of a Prime Number]
\label{thm:totient_prime}
The Totient of a Prime Number: If $m \in \mathbb{N}$ so that $m$ is prime, then $\varphi(m) = m - 1$.
\end{theorem}

\begin{center}
    \vspace{10pt}
    \line(1,0){450}
    \vspace{25pt}
\end{center}

\begin{theorem}[Primality Test via the Totient Function]
\label{thm:primality_test}
A Primality Test via the Totient Function: $\forall m \in \mathbb{N} : m \geq 2$, $\varphi(m) = m - 1 \iff m$ is a prime number.
\end{theorem}
\begin{proof}
    \leavevmode
\begin{itemize}
    \item[($\Rightarrow$)] \textbf{(``Only if'' direction)}: We prove the contrapositive. Assume $m$ is composite (not prime). $\implies \exists p \in \mathbb{N} : 1 < p < m \text{ and } p \mid m$.
    Since $p \mid m$, $\gcd(p, m) = p > 1$. Thus, $p$ is not coprime to $m$. Since $p \in \{1, \dots, m-1\}$ is excluded from the count, the total count $\varphi(m)$ must be strictly less than $m-1$. More precisely, $\varphi(m) \leq m - 2$. This shows that the equality $\varphi(m) = m - 1$ holds $\implies m$ is prime.
    \item[($\Leftarrow$)] \textbf{(``If'' direction)}: Assume $m$ is prime. Then $\forall a \in \{1, \dots, m-1\}$, $\gcd(a, m) = 1$. Since there are $m-1$ such integers, $\varphi(m) = m - 1$. \qedhere
\end{itemize}
\end{proof}

\begin{center}
    \vspace{10pt}
    \line(1,0){450}
    \vspace{25pt}
\end{center}

\begin{proposition}[Totient of a Prime Power]
\label{prop:prime_power_formula}
$\forall p$ prime and $\forall \alpha \in \mathbb{N} : \alpha \geq 1$:
$$ \varphi(p^\alpha) = p^\alpha - p^{\alpha-1} $$
\end{proposition}
\begin{proof}
The totient function $\varphi(p^\alpha)$ counts the number of integers $a \in \{1, \dots, p^\alpha\}$ so that $\gcd(a, p^\alpha) = 1$.
The integers that are \textbf{not} coprime to $p^\alpha$ are those $a$ for which $p \mid a$. We count these integers (the complement).
These non-coprime integers are the multiples of $p$ in the range $[1, p^\alpha]$, i.e., of the form $k \cdot p$ where $k \in \mathbb{N}$.
The condition $1 \leq k \cdot p \leq p^\alpha$ implies $1 \leq k \leq p^{\alpha-1}$.
Thus, there are exactly $p^{\alpha-1}$ multiples of $p$ in the specified range.
Subtracting this count from the total number of integers, $p^\alpha$, yields:
$$ \varphi(p^\alpha) = p^\alpha - p^{\alpha-1} = p^{\alpha-1} (p - 1) $$ \qedhere
\end{proof}

\begin{center}
    \vspace{10pt}
    \line(1,0){450}
    \vspace{25pt}
\end{center}

\begin{theorem}[Multiplicativity of Euler's Totient Function]
\label{thm:multiplicativity}
Multiplicativity of Euler's Totient Function: If $m, n \in \mathbb{N}$ are coprime integers ($\gcd(m, n) = 1$), then $\varphi(m n) = \varphi(m) \varphi(n)$.
\end{theorem}

\newpage

%====================Lesson 4=========================%

\newpage
\section{The fast powering algorithm}

\subsection{Theorem: Existence of Solutions for a System of Two Linear Congruences}

\begin{theorem}[Existence of Solutions for Two Linear Congruences]
Let $m$ and $n$ be positive integers, and let $a$ and $b$ be any integers. The system of linear congruences:
$$
\begin{cases}
x \equiv a \pmod{m} \\
x \equiv b \pmod{n}
\end{cases}
$$
has a solution if and only if $\gcd(m, n) \mid (a - b)$.
\end{theorem}
\begin{proof}[ 1: Constructive Proof by Cases (Alessandra's Method)]
We prove the biconditional nature by demonstrating necessity ($\implies$) and sufficiency ($\impliedby$) separately.

\textbf{1. Necessity ($\implies$)}: We assume a solution exists $\rightarrow$ $\exists x_0 \in \mathbb{Z}$ : $x_0 \equiv a \pmod{m}$ and $x_0 \equiv b \pmod{n}$.
This implies:
$$
x_0 - a = mk \quad \text{for some } k \in \mathbb{Z}
$$
$$
x_0 - b = nl \quad \text{for some } l \in \mathbb{Z}
$$
Thus, $x_0 = a + mk$ and $x_0 = b + nl$. Equating the two expressions for $x_0$ yields $a + mk = b + nl$.
Rearranging the terms: $a - b = nl - mk$.
Let $d = \gcd(m, n)$. By definition, $d \mid m$ and $d \mid n$. Therefore, $d$ must divide any integer linear combination of $m$ and $n$, $\forall k, l \in \mathbb{Z}$.
Since $a - b = nl - mk$, we conclude that $d \mid (a - b)$.

\textbf{2. Sufficiency ($\impliedby$)}: We assume $\gcd(m, n) \mid (a - b)$.
Let $d = \gcd(m, n)$. Since $d \mid (a - b)$, we have $a - b = kd$ for some integer $k \in \mathbb{Z}$.
By Bezout's identity, $\exists u, v \in \mathbb{Z}$ : $d = mu + nv$.
Substituting the expression for $d$: $a - b = k(mu + nv) \rightarrow a - b = kmu + knv$.
Rearranging the equation: $a - kmu = b + knv$.
Let $x_0 = a - kmu$. By construction:
$x_0 = a - kmu \rightarrow x_0 - a = (-ku)m$. Since $-ku \in \mathbb{Z}$, $x_0 \equiv a \pmod{m}$.
Also, $x_0 = b + knv \rightarrow x_0 - b = (kv)n$. Since $kv \in \mathbb{Z}$, $x_0 \equiv b \pmod{n}$.
Since $x_0$ satisfies both congruences, a solution exists.
\end{proof}

\begin{proof}[2: Direct Proof via Reduction (Professor's Method)]
The first congruence $x \equiv a \pmod{m}$ can be rewritten as $x = a + km$ for some integer $k \in \mathbb{Z}$.
Substituting this into the second congruence yields:
$$
a + km \equiv b \pmod{n}
$$
Rearranging the terms, we obtain a single linear congruence for the unknown $k$:
$$
km \equiv b - a \pmod{n}
$$
A solution for this linear congruence exists if and only if $\gcd(m, n) \mid (b - a)$.
Since $\gcd(m, n) = \gcd(m, n)$ and $\gcd(m, n) \mid (b - a)$ is equivalent to $\gcd(m, n) \mid (a - b)$, the original system has a solution if and only if the condition holds.
\end{proof}

\begin{center}
    \line(1,0){450}
\end{center}

\begin{proposition}[Generalization for $n$-Congruence Systems]
A system of congruences $x \equiv a_i \pmod{m_i}$ for $i \in \{1, 2, \ldots, n\}$ has a solution if and only if for every pair of indices $i \neq j$, the condition $\gcd(m_i, m_j) \mid (a_j - a_i)$ holds.
\end{proposition}

\bigskip

\subsection{The Fast Powering Algorithm}
\begin{remark}[Computational Problem]
The algorithm is designed to efficiently calculate the value of $g^A \pmod{N}$, given integers $g > 1, A > 1,$ and $N > 1$. The challenge arises when the exponent $A$ is very large, $\forall A \in \mathbb{N} : A \gg 1$. The naive method requires $\mathcal{O}(A)$ multiplications, which is computationally infeasible for large $A$.
\end{remark}

\begin{remark}[: Fast Powering Method Principle]
The Fast Powering Algorithm leverages the binary representation of the exponent $A$. Let the binary expansion be $A = \sum_{i=0}^r A_i 2^i$, where $A_i \in \{0, 1\}$.
The exponentiation property allows us to write:
$$
g^A = g^{\sum_{i=0}^r A_i 2^i} = \prod_{i=0}^r (g^{2^i})^{A_i} \pmod{N}
$$
The core idea is **Successive Squaring**: the terms $g^{2^i} \pmod{N}$ are computed efficiently by starting with $g^1$ and repeatedly squaring the previous result, $\forall i \in \{1, \ldots, r\}$ : $g^{2^i} = (g^{2^{i-1}})^2 \pmod{N}$.
The final result is the product of only those terms $g^{2^i} \pmod{N}$ for which the corresponding binary coefficient $A_i = 1$.
\end{remark}

\begin{remark}[Re 2.4: Complexity Analysis]
The Fast Powering Algorithm significantly reduces complexity from $\mathcal{O}(A)$ (naive method) to $\mathcal{O}(\log_2(A))$.
This is due to the three phases:
\begin{itemize}
    \item Binary Conversion: $\sim \log_2(A)$ divisions.
    \item Successive Squaring: $\sim \log_2(A)$ squarings.
    \item Final Product: at most $\sim \log_2(A)$ multiplications.
\end{itemize}
The total number of operations $\sim 3 \log_2(A)$, providing immense efficiency gain, $\forall A \in \mathbb{N} : A \approx 2^{1000} \rightarrow \text{Operations} \approx 3000$.
\end{remark}

\begin{center}
    \line(1,0){450}
\end{center}

\subsection{Fermat's Little Theorem}
\begin{theorem}[Th 3.1: Fermat's Little Theorem - Primary Form]
Let $p$ be a prime number. For any integer $a \in \mathbb{Z}$ so that $p$ does not divide $a$, we have:
$$
a^{p-1} \equiv 1 \pmod{p}
$$
\end{theorem}

\begin{theorem}[Th 3.1: Fermat's Little Theorem - Alternative Form]
Let $p$ be a prime number. For any integer $a \in \mathbb{Z}$, we have:
$$
a^p \equiv a \pmod{p}
$$
\end{theorem}

\begin{remark}[Re 3.2: Conditions for Application]
The theorem is generally false for composite moduli. The primary condition is that the modulus $p$ must be a prime number.
\end{remark}


\newpage
\section{A Compendium of Theorems in Modular Arithmetic}

\subsection{Fermat's Little Theorem and Its Applications}

\begin{theorem}[Fermat's Little Theorem]
Let $p$ be a prime. If $a \in (\mathbb{Z}/p\mathbb{Z})^*$, then $a^{p-1} = 1$.
Equivalently, if $p$ is a prime and $a$ is an integer so that $p$ does not divide $a$, then $a^{p-1} \equiv 1 \pmod{p}$.
\end{theorem}
\begin{proof}
The proof employs a technique involving the set of invertible elements modulo $p$.

\textbf{Step 1: Construct the Set}
Let $a \in (\mathbb{Z}/p\mathbb{Z})^*$. Consider the set $S$ formed by multiplying $a$ by each of the $p-1$ invertible elements of $\mathbb{Z}/p\mathbb{Z}$, which are $I = \{1, 2, \ldots, p-1\}$.
$$
S = \{ a \cdot i \pmod{p} \mid \forall i \in I \}
$$

\textbf{Step 2: Prove Non-Zero Property}
We claim that $a \cdot i \not\equiv 0 \pmod{p}$ for all $i \in I$.
Assume for contradiction that $a \cdot i \equiv 0 \pmod{p}$. This implies $p \mid (a \cdot i)$. Since $p$ is prime, $p \mid a$ or $p \mid i$.
By premise, $a \in (\mathbb{Z}/p\mathbb{Z})^*$, so $p \nmid a$.
Since $1 \le i \le p-1$, $\gcd(i, p) = 1$, so $p \nmid i$.
This is a contradiction, so $\forall i \in I$, $a \cdot i \not\equiv 0 \pmod{p}$.

\textbf{Step 3: Prove Distinctness}
We claim that all $p-1$ elements in $S$ are distinct.
Assume for contradiction that $\exists i, j \in I$ so that $i \ne j$ and $a \cdot i \equiv a \cdot j \pmod{p}$.
This implies $a \cdot (i-j) \equiv 0 \pmod{p}$. Since $p \nmid a$, we must have $p \mid (i-j)$.
However, since $i, j \in \{1, \ldots, p-1\}$ and $i \ne j$, the difference $i-j$ satisfies $0 < |i-j| < p$. Such an integer cannot be a multiple of $p$.
This contradiction proves that all elements in $S$ are distinct.

\textbf{Step 4: Equate Sets}
Since $S$ contains $p-1$ distinct, non-zero elements modulo $p$, and the set of all non-zero elements is $(\mathbb{Z}/p\mathbb{Z})^* = I = \{1, 2, \ldots, p-1\}$, the two sets must be identical, $\forall i \in I : S = I$.

\textbf{Step 5: The Trick (Multiplication)}
The product of the elements in $S$ must equal the product of the elements in $I$:
$$
\prod_{i=1}^{p-1} (a \cdot i) \equiv \prod_{i=1}^{p-1} i \pmod{p}
$$

\textbf{Step 6: Simplify and Conclude}
Factoring the left-hand side yields:
$$
a^{p-1} \cdot (p-1)! \equiv (p-1)! \pmod{p}
$$
Since $p$ is prime, $(p-1)!$ is invertible modulo $p$ (by Wilson's Theorem, or simply because $\forall i \in \{1, \ldots, p-1\}$, $\gcd(i, p)=1$). We can safely cancel $(p-1)!$ from both sides:
$$
a^{p-1} \equiv 1 \pmod{p}
$$
\end{proof}

\begin{center}
    \line(1,0){450}
\end{center}

\begin{corollary}[Modular Exponentiation]
Let $p$ be a prime, and $a \in \mathbb{Z}$. If $\forall x, y \in \mathbb{N} \setminus \{0\}$ so that $x \equiv y \pmod{p-1}$, then $a^x \equiv a^y \pmod{p}$.
\end{corollary}
\begin{proof}
We consider two distinct cases for the integer $a$.

\textbf{Case 1: $p$ divides $a$}
If $p \mid a$, then $a \equiv 0 \pmod{p}$. Since $x \ge 1$ and $y \ge 1$, we have $a^x \equiv 0 \pmod{p}$ and $a^y \equiv 0 \pmod{p}$.
Thus, $a^x \equiv a^y \pmod{p}$ holds trivially.

\textbf{Case 2: $p$ does not divide $a$}
The congruence $x \equiv y \pmod{p-1}$ means $y = x + k(p-1)$ for some $k \in \mathbb{Z}$.
Substituting this into $a^y$:
$$
a^y = a^{x + k(p-1)} = a^x \cdot a^{k(p-1)} = a^x \cdot (a^{p-1})^k
$$
By Fermat's Little Theorem, $a^{p-1} \equiv 1 \pmod{p}$. Therefore:
$$
a^y \equiv a^x \cdot (1)^k \pmod{p} \rightarrow a^y \equiv a^x \pmod{p}
$$
\end{proof}

\begin{center}
    \line(1,0){450}
\end{center}

\begin{corollary}[Finding Modular Inverses]
Let $p$ be a prime and $a \in (\mathbb{Z}/p\mathbb{Z})^*$. Then the inverse of $a$ modulo $p$ is $a^{p-2}$.
\end{corollary}
\begin{proof}
We must show that $a \cdot a^{p-2} \equiv 1 \pmod{p}$.
$$
a \cdot a^{p-2} = a^{1 + p-2} = a^{p-1}
$$
By Fermat's Little Theorem, $\forall a \in \mathbb{Z}$ so that $p \nmid a$, we have $a^{p-1} \equiv 1 \pmod{p}$.
Therefore:
$$
a \cdot a^{p-2} \equiv 1 \pmod{p}
$$
This confirms that $a^{p-2}$ is the multiplicative inverse of $a$ modulo $p$.
\end{proof}

\begin{center}
    \line(1,0){450}
\end{center}

\begin{corollary}[Fermat Primality Test]
Let $p$ be an integer so that $p \ge 2$, and let $a$ be an integer so that $p$ does not divide $a$. If $a^{p-1} \not\equiv 1 \pmod{p}$, then $p$ is not a prime number.
\end{corollary}
\begin{proof}
This statement is the contrapositive of Fermat's Little Theorem:
$(\neg \text{THEN}) \implies (\neg \text{IF})$.
The original theorem states: $\text{IF } p \text{ is prime, THEN } a^{p-1} \equiv 1 \pmod{p}$.
The contrapositive is: $\text{IF } a^{p-1} \not\equiv 1 \pmod{p}, \text{ THEN } p \text{ is not prime}$.
This conclusion is logically sound.
\end{proof}

\bigskip



%====================Lesson 4=========================%

\newpage
\section{The Order of an Element and Primitive Roots}

\begin{definition}[Order of an Element]
Let $p$ be a prime and $a \in \mathbb{Z}$ so that $p$ does not divide $a$. The order of $a$ modulo $p$, denoted $\text{ord}_p(a)$, is the minimum positive integer $n \in \mathbb{N} \setminus \{0\}$ so that $a^n \equiv 1 \pmod{p}$.
\end{definition}

\begin{center}
    \vspace{10pt}
    \line(1,0){450}
    \vspace{25pt}
\end{center}

\begin{proposition}[Order Divisibility Property]
Let $p$ be a prime, $a \in \mathbb{Z}$ so that $p$ does not divide $a$, and $n \ge 1$. Then $a^n \equiv 1 \pmod{p}$ if and only if $\text{ord}_p(a)$ divides $n$.
\end{proposition}
\begin{proof}
We demonstrate both directions of the biconditional.

\textbf{First Direction ($\impliedby$)}: Assume $\text{ord}_p(a) \mid n$.
$\rightarrow \exists k \in \mathbb{Z}$ so that $n = k \cdot \text{ord}_p(a)$.
Then:
$$
a^n = a^{k \cdot \text{ord}_p(a)} = (a^{\text{ord}_p(a)})^k
$$
By definition of the order, $a^{\text{ord}_p(a)} \equiv 1 \pmod{p}$. Therefore:
$$
a^n \equiv (1)^k \equiv 1 \pmod{p}
$$

\textbf{Second Direction ($\implies$)}: Assume $a^n \equiv 1 \pmod{p}$.
By the division algorithm, $\exists q, r \in \mathbb{Z}$ so that $n = q \cdot \text{ord}_p(a) + r$, where $0 \le r < \text{ord}_p(a)$.
We evaluate $a^n$:
$$
a^n = a^{q \cdot \text{ord}_p(a) + r} = (a^{\text{ord}_p(a)})^q \cdot a^r
$$
Substituting the assumed and known congruences:
$$
1 \equiv (1)^q \cdot a^r \pmod{p} \rightarrow a^r \equiv 1 \pmod{p}
$$
Since $\text{ord}_p(a)$ is the *minimum* positive integer for which this holds, and $0 \le r < \text{ord}_p(a)$, the only possibility is that $r = 0$.
$\rightarrow \text{ord}_p(a)$ divides $n$.
\end{proof}

\begin{center}
    \vspace{10pt}
    \line(1,0){450}
    \vspace{25pt}
\end{center}

\begin{proposition}[Equality of Powers]
Let $p$ be a prime and $g \in (\mathbb{Z}/p\mathbb{Z})^*$. Then $g^x = g^y$ if and only if $x \equiv y \pmod{\text{ord}_p(g)}$.
\end{proposition}
\begin{proof}
Since $g \in (\mathbb{Z}/p\mathbb{Z})^*$, $g$ is invertible.
$$
g^x = g^y \iff g^x \cdot g^{-y} = 1 \iff g^{x-y} = 1
$$
Applying the proposition from the previous section (2.1), $g^n = 1$ if and only if $\text{ord}_p(g) \mid n$.
Let $n = x-y$. We conclude:
$$
g^{x-y} = 1 \iff \text{ord}_p(g) \mid (x-y)
$$
By the definition of modular congruence, this is equivalent to: $x \equiv y \pmod{\text{ord}_p(g)}$.
\end{proof}

\begin{center}
    \vspace{10pt}
    \line(1,0){450}
    \vspace{25pt}
\end{center}

\begin{definition}[Primitive Root]
An element $g \in (\mathbb{Z}/p\mathbb{Z})^*$ is a primitive root if its powers generate all the elements of $(\mathbb{Z}/p\mathbb{Z})^*$. That is, the set $\{ g^1, g^2, \ldots, g^{p-1} \}$ is equal to the set $\{ 1, 2, \ldots, p-1 \}$.
\end{definition}

\begin{center}
    \vspace{10pt}
    \line(1,0){450}
    \vspace{25pt}
\end{center}

\begin{theorem}[Primitive Root Condition]
An element $g \in (\mathbb{Z}/p\mathbb{Z})^*$ is a primitive root if and only if $\text{ord}_p(g) = p-1$.
\end{theorem}
\begin{proof}[Proof Outline]
For $g$ to be a primitive root, the $p-1$ powers $\{ g^1, g^2, \ldots, g^{p-1} \}$ must all be distinct.
By the Proposition on the Condition for Equality of Powers, $g^i = g^j$ if and only if $i \equiv j \pmod{\text{ord}_p(g)}$.
For the powers to be distinct for all $1 \le i < j \le p-1$, we must ensure that $g^{j-i} \ne 1$ for any exponent $j-i$ where $1 \le j-i < p-1$.
This means that the smallest positive exponent $k$ so that $g^k = 1$ must be $p-1$.
By definition, $\text{ord}_p(g) = k$.
Therefore, $g$ is a primitive root if and only if $\text{ord}_p(g) = p-1$.
\end{proof}

\bigskip
\begin{definition}
    Let $p$ be a prime, and $a \in \mathbb{Z}$ such that $p \nmid a$.
    The \textbf{order} $ord_{p}(a)$ of $a$ modulo $p$ is the minimum $n \in \mathbb{N}$ such that $n \geq 1$ and $a^{n} \equiv 1 \pmod{p}$.
\end{definition}

\bigskip
\begin{proposition}
    Let $p$ be a prime, $a \in \mathbb{Z}$ such that $p \nmid a$, and $n \ge 1$.
    Then $a^{n} \equiv 1 \pmod{p}$ if and only if the order of $a$ modulo $p$, denoted $\text{ord}_{p}(a)$, divides $n$.
    In particular, $\text{ord}_{p}(a) \mid (p-1)$ (a consequence of Fermat's Little Theorem).
\end{proposition}

\begin{proof}
    Let $k = \text{ord}_{p}(a)$.
    
    \textbf{$(\Rightarrow)$ If $a^{n} \equiv 1 \pmod{p}$, then $k \mid n$} \\
    Use the Division Algorithm on $n$ by $k$: $n = qk + r$, where $0 \le r < k$.
    
    $$1 \equiv a^{n} \equiv a^{qk + r} \equiv (a^{k})^{q} \cdot a^{r} \equiv 1^{q} \cdot a^{r} \equiv a^{r} \pmod{p}$$
    
    Since $a^{r} \equiv 1 \pmod{p}$ and $k$ is the \textbf{minimum positive exponent} satisfying this property, and $r < k$, we must have $r=0$.
    Thus, $n = qk$, which means $k \mid n$.

    \textbf{$(\Leftarrow)$ If $k \mid n$, then $a^{n} \equiv 1 \pmod{p}$} \\
    If $k \mid n$, then $n = qk$ for some $q \in \mathbb{Z}$.
    $$a^{n} = a^{qk} = (a^{k})^{q} \equiv 1^{q} \equiv 1 \pmod{p}$$
\end{proof}

\begin{center}
    \vspace{10pt}
    \line(1,0){450}
    \vspace{25pt}
\end{center}

\begin{corollary}
    Let $p$ be a prime, $a \in \mathbb{Z}$ such that $p \nmid a.$
    $ord_{p}(a) = \min\{n \in \mathbb{N} : n \geq 1 \text{ so that } a^{n} \equiv 1 \pmod{p}\}$
    $= \min\{n \in \mathbb{N} : n \geq 1 \text{ so that } a^{n} \equiv 1 \pmod{p} \text{ and } n \mid p-1\}$.
\end{corollary}

\begin{center}
    \vspace{10pt}
    \line(1,0){450}
    \vspace{25pt}
\end{center}

\begin{proposition}
    Let $p$ be a prime. If $g \in \mathbb{F}_{p}^{*}$, then $g^{x} = g^{y} \iff y \equiv x \pmod{ord_{p}(g)}$.
\end{proposition}
\begin{proof}
    ($\Leftarrow$) Assume $x \equiv y \pmod{ord_{p}(g)}$. Then $x-y = k \cdot ord_{p}(g)$ for some integer $k$.
    $g^{x} = g^{y+k \cdot ord_{p}(g)} = g^{y} \cdot (g^{ord_{p}(g)})^{k} \equiv g^{y} \cdot 1^{k} \equiv g^{y} \pmod{p}$.

    ($\Rightarrow$) Assume $g^{x} = g^{y}$. Without loss of generality, assume $x \geq y$.
    Since $g \in \mathbb{F}_{p}^{*}$ is invertible, multiplying by $(g^{-1})^{y}$ gives $g^{x-y} = 1$.
    By the previous Proposition, $ord_{p}(g) \mid x-y$, which means $x \equiv y \pmod{ord_{p}(g)}$.
\end{proof}

\begin{center}
    \vspace{10pt}
    \line(1,0){450}
    \vspace{25pt}
\end{center}

\begin{definition}
    Let $p$ be a prime.
    A \textbf{primitive root} of $\mathbb{F}_{p}$ is any $g \in \mathbb{F}_{p}^{*}$ such that the powers of $g$ generate all the elements of $\mathbb{F}_{p}^{*}$, i.e., $\mathbb{F}_{p}^{*} = \{1, g, g^{2}, \dots, g^{p-2}\}$.
\end{definition}

\begin{center}
    \vspace{10pt}
    \line(1,0){450}
    \vspace{25pt}
\end{center}

\begin{theorem}[Characterization of the Primitive Roots]
    Let $p$ be a prime. $g \in \mathbb{F}_{p}^{*}$ is a primitive root if and only if $ord_{p}(g) = p-1$.
\end{theorem}
\begin{proof}
    The proof establishes the equivalence between $g$ generating the group $\mathbb{F}_{p}^{*}$ and $g$ having the maximum possible order. Note that $|\mathbb{F}_{p}^{*}| = p-1$.

    \textbf{$(\Rightarrow)$ If $g$ is a primitive root, then $\text{ord}_{p}(g) = p-1$} \\
    By definition, if $g$ is a primitive root, the set of its powers $\{1, g, \dots, g^{p-2}\}$ contains \textbf{all $p-1$ distinct elements} of $\mathbb{F}_{p}^{*}$.
    
    Let $m = \text{ord}_{p}(g)$. If we assume $m < p-1$, then the sequence of powers $1, g, g^2, \dots$ repeats with a period of $m$. This means the set of distinct powers $\langle g \rangle$ has size $|\langle g \rangle| = m$.
    
    Since $g$ is a primitive root, $\langle g \rangle = \mathbb{F}_{p}^{*}$, so $m = |\mathbb{F}_{p}^{*}| = p-1$.
    The assumption $m < p-1$ leads to the contradiction $|\langle g \rangle| < |\mathbb{F}_{p}^{*}|$. Thus, $m = p-1$.

    \textbf{$(\Leftarrow)$ If $\text{ord}_{p}(g) = p-1$, then $g$ is a primitive root} \\
    Assume $\text{ord}_{p}(g) = p-1$. We must show the powers $1, g, \dots, g^{p-2}$ are all distinct.
    
    Suppose, for contradiction, that $g^{i} \equiv g^{j} \pmod{p}$ for $0 \le i < j \le p-2$.
    Since $g$ is invertible, $g^{j-i} \equiv 1 \pmod{p}$.
    
    Let $k = j-i$. Then $0 < k < p-1$.
    This implies that a power $g^k \equiv 1 \pmod{p}$ for an exponent $k$ smaller than $p-1$. This contradicts the assumption that $\text{ord}_{p}(g) = p-1$ is the \textbf{minimum} positive exponent that results in $1$.
    
    Thus, the $p-1$ powers $\{1, g, \dots, g^{p-2}\}$ are distinct. Since $\mathbb{F}_{p}^{*}$ has exactly $p-1$ elements, the powers must generate the entire group, meaning $g$ is a primitive root.
\end{proof}

\begin{center}
    \line(1,0){450}
\end{center}

\begin{theorem}[Existence of Primitive Roots]
    Let $p$ be a prime. $\mathbb{F}_{p}$ has at least one primitive root. (Admitted)
\end{theorem}

\bigskip

\begin{proof}
    Let $x = g^{\frac{p-1}{2}}$. We first observe the square of $x$:
    $$x^{2} = \left(g^{\frac{p-1}{2}}\right)^{2} = g^{p-1}$$
    
    By \textbf{Fermat's Little Theorem}, $g^{p-1} \equiv 1 \pmod{p}$, provided $p \nmid g$ (which is true since $g \in \mathbb{F}_p^*$).
    Thus, $x$ satisfies the congruence:
    $$x^{2} \equiv 1 \pmod{p}$$
    
    In the field $\mathbb{F}_{p}$, the only solutions to $x^{2} \equiv 1$ are $x=1$ and $x=-1$ (or $x \equiv p-1 \pmod{p}$).
    
    We must show that $x \neq 1$. We proceed by contradiction.
    
    Assume $x=1$, meaning:
    $$g^{\frac{p-1}{2}} \equiv 1 \pmod{p}$$
    
    By the property of the order of an element, this would imply:
    $$\text{ord}_{p}(g) \mid \frac{p-1}{2}$$
    
    However, since $g$ is a \textbf{primitive root}, we know by definition that $\text{ord}_{p}(g) = p-1$.
    The divisibility condition becomes:
    $$(p-1) \mid \frac{p-1}{2}$$
    
    Since $p>2$, this condition is false (a positive integer cannot divide a number half its size). This is a \textbf{contradiction}.
    
    Therefore, the only remaining possibility for the value of $x$ is $x=-1$.
    $$g^{\frac{p-1}{2}} \equiv -1 \pmod{p}$$
\end{proof}

\begin{center}
    \vspace{10pt}
    \line(1,0){450}
    \vspace{25pt}
\end{center}

\begin{remark}
    Let $p>2$ be a prime. If $g$ is a square in $\mathbb{F}_{p}$, then $g$ is not a primitive root in $\mathbb{F}_{p}$.
\end{remark}
\begin{proof}
    Let $g$ be a square, so $g = h^{2}$ for some $h \in \mathbb{F}_{p}^{*}$.
    Then $g^{(p-1)/2} = (h^{2})^{(p-1)/2} = h^{p-1} \equiv 1 \pmod{p}$ by Fermat's Little Theorem.
    Since $g^{(p-1)/2} \equiv 1 \pmod{p}$, by the proposition on order, $ord_{p}(g) \mid \frac{p-1}{2}$.
    Since $\frac{p-1}{2} < p-1$ (as $p>2$), we have $ord_{p}(g) < p-1$.
    By the characterization theorem, $g$ is not a primitive root.
\end{proof}

\begin{center}
    \vspace{10pt}
    \line(1,0){450}
    \vspace{25pt}
\end{center}

\begin{lemma}[Effective Test for a Primitive Root Candidate]
    Let $p>2$ be prime and let $p-1 = \prod_{i=1}^{m}p_{i}^{a_{i}}$ be the prime factorization of $p-1$.
    Then $g \in \mathbb{F}_{p}^{*}$ is a primitive root if and only if $g^{(p-1)/p_{i}} \not\equiv 1 \pmod{p}$ for all $i=1, \dots, m$.
\end{lemma}

\begin{proof}
    Let $k = \text{ord}_{p}(g)$. Recall that $g$ is a primitive root $\iff k = p-1$.
    
    \textbf{$(\Rightarrow)$ If $g$ is a primitive root, the test passes} \\
    Assume $g$ is a primitive root, so $k = p-1$.
    
    Suppose, for contradiction, that $g^{(p-1)/p_{i}} \equiv 1 \pmod{p}$ for some prime factor $p_i$.
    By the property of the order, $k$ must divide the exponent:
    $$k \mid \frac{p-1}{p_{i}} \implies (p-1) \mid \frac{p-1}{p_{i}}$$
    Since $p_i$ is a prime ($p_i \ge 2$), this divisibility is impossible.
    Thus, $g^{(p-1)/p_{i}} \not\equiv 1 \pmod{p}$ must hold for all $i$.

    \textbf{$(\Leftarrow)$ If the test passes, $g$ is a primitive root} \\
    Assume $g^{(p-1)/p_{i}} \not\equiv 1 \pmod{p}$ for all $i$.
    We know $k \mid p-1$, so $k$ is a divisor of $p-1$, and thus $k$ must be of the form $k = p_{1}^{f_{1}} \cdots p_{m}^{f_{m}}$ with $0 \le f_{i} \le a_{i}$.
    
    Suppose, for contradiction, that $k < p-1$. This means at least one exponent $f_i$ must be strictly less than $a_i$ (i.e., $f_i \le a_i - 1$).
    
    For such an index $i$, we can construct the exponent $\frac{p-1}{p_{i}}$. Since $p_i^{a_i}$ is a factor of $p-1$ and $p_i^{f_i}$ is a factor of $k$, the difference in the exponent is:
    $$\frac{p-1}{p_{i}} = \frac{p_{1}^{a_{1}} \cdots p_{i}^{a_{i}-1} \cdots p_{m}^{a_{m}}}{p_{i}}$$
    
    Since $f_i \le a_i - 1$, we have $k \mid \frac{p-1}{p_{i}}$.
    
    By the property of the order, $k \mid \frac{p-1}{p_{i}}$ implies $g^{(p-1)/p_{i}} \equiv 1 \pmod{p}$.
    This contradicts our initial assumption that $g^{(p-1)/p_{i}} \not\equiv 1 \pmod{p}$ for all $i$.
    
    Therefore, the assumption $k < p-1$ is false, meaning $k = p-1$. Hence, $g$ is a primitive root. \qedhere
\end{proof}

\begin{center}
    \vspace{10pt}
    \line(1,0){450}
    \vspace{25pt}
\end{center}

\begin{algorithm}[Finding a Primitive Root]
    Let $p>2$ be prime and $p-1 = \prod_{i=1}^{m}p_{i}^{a_{i}}$ be the prime factorization of $p-1$, with $p_{i}$ distinct primes and $a_{i} \geq 1$ for $i=1, \dots, m$.
    \begin{enumerate}
        \item Start with $g=2$.
        \item If $g^{(p-1)/p_{i}} \not\equiv 1 \pmod{p}$ for all $i=1, \dots, m$, then $g$ is a primitive root.
        \item If not, increment $g \rightarrow g+1$ and return to Step 2.
    \end{enumerate}
\end{algorithm}


\end{document}